
\SGARefTarget{sga3:xxiv:1.12}
\paragraph{\SGARefLink{sga3:xxiv:1.12}{Corollary 1.12}.}
\emph{Let \(S\) be a henselian local scheme, let \(G\) be a
reductive \(S\)-group, let \(G'\) be a smooth affine \(S\)-group with
connected fibres, and let \(s\) be the closed point of \(S\). If
\(G_s\) and \(G'_s\) are isomorphic algebraic groups over
\(\kappa(s)\), then \(G\) and \(G'\) are isomorphic. More precisely,
every \(\kappa(s)\)-isomorphism}
\[
  G_s\simeq G'_s
\]
\emph{comes from an \(S\)-isomorphism \(G\simeq G'\).}

Applying \SGARefLink{sga3:xxiv:1.7}{Corollary~1.7} \textup{(respectively, \SGARefLink{sga3:xxiv:1.12}{Corollary~1.12})} to the
scheme of dual numbers over a field, \SGARefLink{sga3:iii:remark:2-10}{Exp.~III, 2.10}
\textup{(respectively, \SGARefLink{sga3:xxiv:3.10}{3.10})} yields part \textup{(i)}
\textup{(respectively, part \textup{(ii)})} of the next corollary.

\SGARefTarget{sga3:xxiv:1.13}
\paragraph{\SGARefLink{sga3:xxiv:1.13}{Corollary 1.13}.}
\emph{Let \(k\) be a field and \(G\) a reductive \(k\)-group.}
\begin{enumerate}
\item[\SGAInlineRefTarget{sga3:xxiv:1.13:item:i}(i)] \emph{If \(G\) is adjoint, then}
\[
  H^1(G,\Lie(G/k))=0.
\]
\footnote{The source's editorial note corrects the original, which
stated \textup{\SGARefLink{sga3:xxiv:1.13:item:i}{(i)}} without assuming \(G\) adjoint. By the
description of \(H^i(G,-)\) as the derived functors of
\(\operatorname{Hom}_G(k,-)\) \textup{(\SGARefLink{sga3:i:theorem:5-3-1}{Exp.~I, 5.3.1}; see also
[Ja87], I, 4.16)}, one has
\(H^i(G,V)=\operatorname{Ext}^i_G(k,V)\) for every \(G\)-module
\(V\). If \(\operatorname{char}(k)=p>0\) and
\(G=\operatorname{SL}_{p,k}\), then
\(\Lie(\operatorname{GL}_p/k)\) is a nontrivial extension of \(k\)
by \(\mathfrak g=\Lie(\operatorname{SL}_p/k)\), so
\(H^1(G,\mathfrak g)\neq0\). See also the addition \SGARefLink{sga3:xxiv:1.15.1}{1.15.1} below.}
\item[\SGAInlineRefTarget{sga3:xxiv:1.13:item:ii}(ii)] \emph{One has}
\[
  H^2(G,\Lie(G/k))=0.
\]
\end{enumerate}

\SGARefTarget{sga3:xxiv:1.14}
\paragraph{\SGARefLink{sga3:xxiv:1.14}{Remark 1.14}.}
\begin{enumerate}
\item[(i)] The assertion concerning \(H^1\) was known to Chevalley;
the assertion concerning \(H^2\) had been proved in most cases of
the classification by Chevalley.
\item[(ii)] In fact, the conjunction of \SGARefLink{sga3:xxiv:1.13}{Corollary~1.13} and the
uniqueness theorem over an algebraically closed field is essentially
equivalent to the uniqueness theorem. A direct proof of
\SGARefLink{sga3:xxiv:1.13}{Corollary~1.13} would therefore give a way to deduce the general
uniqueness theorem from Chevalley's uniqueness theorem over a field.
\end{enumerate}

The existence of reductive groups of every type over every scheme
\textup{(Exp.~XXV)} shows that the obstructions to lifting a
reductive \(k\)-group \(G\) over artinian rings with residue field
\(k\), which by \SGARefLink{sga3:iii:corollary:3-8}{Exp.~III, 3.8} are elements of
\[
  H^3(G,\Lie(G/k)\otimes V)
  \simeq H^3(G,\Lie(G/k))\otimes V,
\]
where \(V\) is a certain \(k\)-vector space endowed with the trivial
action of \(G\), vanish. This suggests that
\[
  H^3(G,\Lie(G/k))=0.
\]
A direct proof, if the assertion is true, would likewise doubtless
give a way to deduce the general existence theorem from the
existence theorem over a field \textup{(Chevalley's Tôhoku)}.

\SGARefTarget{sga3:xxiv:1.15}
\paragraph{\SGARefLink{sga3:xxiv:1.15}{Corollary 1.15}.}
\emph{Let \(k\) be a field,\footnote{The source's editorial note
replaces ``scheme'' with ``field.''} and let \(G\) be a reductive
\(k\)-group. Regard \(k\) as a trivial \(G\)-module. Then}
\[
  H^1(G,k)=H^2(G,k)=0.
\]
\footnote{Because of the correction made in \SGARefLink{sga3:xxiv:1.13}{1.13}, the source's
editors supply another proof for \(H^1\). Moreover, a theorem of
G.~Kempf gives \(H^i(G,k)=0\) for every \(i>0\); cf.\ [Ja87],
II, \SGARefLink{sga3:ii:proposition:4-5}{4.5} and 4.11.}
Since
\[
  H^i(G_{\bar k},\bar k)=H^i(G,k)\otimes_k\bar k,
\]
we may assume \(k\) algebraically closed. An element of \(H^1(G,k)\)
is simply a morphism of \(k\)-groups
\[
  \varphi:G\longrightarrow\mathbb G_{a,k}.
\]
Then
\[
  \varphi(G)=G/\ker(\varphi)
\]
is a smooth connected reductive subgroup of
\(\mathbb G_{a,k}\) \textup{(cf.\ XIX, \SGARefLink{sga3:xix:paragraph:1-7}{1.7})}, and is therefore
trivial. Thus \(H^1(G,k)=0\).
